% Frozen before running the Aer shot study. Compile with xelatex.
\documentclass[11pt,a4paper]{article}
\usepackage{kotex}
\usepackage{amsmath,amssymb,booktabs,geometry}
\geometry{margin=25mm}
\title{I01 선검증 3 사전 프로토콜: 비균일 표지의 유한 샷 검출력}
\author{RELATIONAL ATLAS}
\date{2026년 10월 4일}
\begin{document}
\maketitle

\section{선택 근거와 제한}
선검증 2의 정확 확률을 사용한 \texttt{qiskit\_rank\_holdout3\_power.py}
파일은 Aer 실행 전에 $n=3,\beta=0.45$, 비선형 시계 에너지,
폭이 다른 단일 복귀 프로젝터를 검정력 시범 조건으로 골랐다.
따라서 이 조건은 새로운 독립 물리 홀드아웃이 아니라 \emph{기존
정확 모형의 유한 샷 구현 검증}이다. 발전연구 주제를 정하지 않는다.

\section{고정 회로와 표본}
$k_i\in\{-1,0,1\}$, $w_k\propto e^{-0.45E_0(k)}$,
$E_0=\sum_i k_i^2$, $E=E_0+0.37E_0^2/3$,
검출 기준 폭 $\beta_r=0.15$, 축별 바이어스 1.25 rad를 사용한다.
$n=3$, $L=7$의 위상 표지 $2\pi t/7$을 주고 공통 과학 위상을
$0\to0.60$ rad 이동한다. 시계 위상 $\tau=0,\pi/2$,
비제약/$P=0$ 상태, 균일/비균일 표지 분포를 모두 실행한다.
균일 $a_t=1/7$, 비균일
$a_t=[1+0.35\cos(2\pi t/7)]/7$이다.

축·시계·상태·위상 블록마다 총 $7\times32768=229376$ shots를
표지에 배분한다. 각 분포의 이론적 $a_t$에 총 shots를 곱해 내림한 뒤
남은 shots를 소수부가 가장 큰 표지 순으로 배분한다. 두 위상 블록은
같은 표지별 shots 수를 갖는다. 표지 기록을 쓴 분석과 표지를
합친 분석은 \emph{동일한 Aer counts}로 한다. CPU workers 6,
Aer statevector, seed 20261017을 고정한다. 회로는 에너지 대각 게이트,
운동량 위상 $R_z$, 기준 상태 역준비 및 0 상태 측정으로 구성한다.

\section{사전 판정}
8개 삭제-표지 두 블록 차이 검정(2 시계$\times$2 상태$\times$2 분포)은
축별 성공수를 합치되 축을 구분한 조건부 재배분 검정 2000회로
$p$값을 계산한다. 가족별 문턱은 Bonferroni $0.05/8=0.00625$다.
정확 모형에서 균일/비제약, 모든 $P=0$는 공통 위상 null이다.
비균일/비제약은 두 시계 위상 모두 정확 양수다.
그러나 파일럿 근사 비중심성은 32768 shots/표지일 때
$\tau=0$에서 약 26.5, $\tau=\pi/2$에서 약 2.1이므로,
두 번째의 미검출은 null의 증거가 아니다. 첫 번째도 한 번의
shot seed에서 반드시 검출된다는 보장은 없다. 사전 문턱과
정확 양수/영을 분리하여 보고한다.

이 실험은 표지별 shots 수를 정확히 알고 배분한 경우다.
무작위 표지 수의 요동이나 미지 분포의 추정 오차는 따로 남는 검증
게이트다. QPU는 사용하지 않는다.
\end{document}
